\section{How Setups Combine and How Context Changes Interpretation}

\subsection{No Setup Lives Alone}

The source material repeatedly rejects one-signal trading.
The trader is not supposed to ask ``is there a density?''
The trader is supposed to ask:
\begin{itemize}
  \item where is the density;
  \item what happened before it;
  \item who is likely trapped around it;
  \item what is the tape doing into it;
  \item what sits behind it;
  \item how stretched the move already is;
  \item and what kind of day this instrument is having.
\end{itemize}

That means the repo should treat setups as \textbf{composite states}, not isolated triggers.

\subsection{A Bounce Can Become a Breakout}

One of the most important transition patterns is that a defense can weaken over time.
A level that initially offers a clean bounce may later become a breakout candidate if:
\begin{itemize}
  \item repeated tests consume the same-side support;
  \item the defending participant refreshes less aggressively;
  \item compression keeps building into the level;
  \item the opposite side becomes more initiative-driven;
  \item nearby support behind the level disappears.
\end{itemize}

Research implication: detector logic should not tag a level once and freeze its role.
The role of a level must be allowed to transition from defended anchor to vulnerable barrier.

\subsection{A Breakout Can Become a False Breakout}

Likewise, a valid breakout trigger can fail if post-break acceptance is weak.
The repository should explicitly model a \emph{post-trigger evaluation window}.
That window needs to answer:
\begin{itemize}
  \item did the market gain distance quickly enough;
  \item did the book behind the break remain supportive;
  \item did the tape continue in the break direction;
  \item did price hold beyond the level or immediately re-enter the prior range.
\end{itemize}

This is not a small nuance.
It is central to the entire source worldview.
Many of the same traders who like breakout trades also like fading failed breakouts.
The transition itself is the opportunity.

\subsection{Static Levels, Dynamic Liquidity, and Executed Flow}

The knowledge base uses three different concepts of ``level'':
\begin{enumerate}
  \item \textbf{static chart level}: previous high, low, range edge, prior extreme;
  \item \textbf{dynamic displayed level}: current density, shifting quote wall, robot layer;
  \item \textbf{executed memory level}: cluster concentration, PoC, prior heavy business.
\end{enumerate}

The best trades often occur when these three notions overlap.
If only one of them exists, quality is usually lower.

\subsection{Confidence and Quality as a Multi-Factor Score}

The transcripts talk about needing several reasons for a trade.
That idea translates naturally into a confidence model.
An eventual research-friendly quality score could include:
\begin{itemize}
  \item structural location quality;
  \item visible-liquidity quality;
  \item tape confirmation quality;
  \item cluster-memory quality;
  \item stretch or spring quality;
  \item room-to-target quality;
  \item day-regime fit;
  \item and execution quality, including spread and queue conditions.
\end{itemize}

This should not initially be over-optimized.
A simple additive score or state-based filter is sufficient for early research.

\subsection{Session Structure and Day-Specific Adaptation}

The materials strongly imply a session-level feedback loop:
\begin{itemize}
  \item observe the early session;
  \item learn whether visible levels are holding or failing today;
  \item learn whether the instrument is moving enough for breakouts;
  \item learn whether fakeouts dominate or whether continuation is clean;
  \item adapt participation accordingly.
\end{itemize}

This is one of the most valuable research directions in the repository because it connects discretionary intuition to systematic adaptation.
Instead of assuming one static detector threshold forever, the system can maintain same-day priors such as:
\begin{itemize}
  \item current breakout follow-through quality;
  \item current bounce hold rate;
  \item current spread and depth regime;
  \item current response to news or leaders;
  \item current instrument liveliness.
\end{itemize}

\subsection{Time-of-Day Conditioning}

The same local pattern should be interpreted differently depending on time of day:
\begin{itemize}
  \item early session often contains the clearest structural information and strongest recalibration opportunity;
  \item midday may have lower conviction or thinner flow;
  \item event windows can create abnormal opportunity and abnormal risk;
  \item late session moves can be distorted by inventory squaring and reduced horizon.
\end{itemize}

For research purposes, time-of-day should not be a cosmetic feature.
It should be a first-class conditioning variable.

\subsection{Clean Decision Checklist}

The human discretionary workflow implied by the materials can be written as a checklist:
\begin{enumerate}
  \item Is the instrument liquid and moving enough right now?
  \item Is there a meaningful structural location?
  \item Is there visible or hidden participant logic at that location?
  \item Does tape or executed flow confirm or contradict the local thesis?
  \item Is the stop precise and behaviorally justified?
  \item Is there room to the next meaningful obstacle?
  \item Does the current day regime fit this setup family?
  \item Is the execution environment acceptable?
\end{enumerate}

If multiple answers are weak, the source material would usually classify the trade as low quality even if a single pattern component is present.
